\subsection{What the decoder comparison can and cannot mean}
\label{sec:probes}

The comparison in the previous subsections is easy to over-read, and the direction of the
over-reading is the flattering one, so we state the ceiling on it explicitly.

\paragraph{The probe class contains the network's own decoder.} The affine probe is fitted on the
post-ReLU state $z=\mathrm{ReLU}(\Phi b)$ with a design matrix that appends an intercept and nothing
else (\code{\_design} in \code{src/lrtr/probes.py}), so its hypothesis class is every affine map of
$z$ --- which includes $W_{\mathrm{out}}$, the map the network's own output layer applies, together
with the network's threshold. It follows that there \emph{exists} a probe configuration attaining the
network's score exactly, for every model in every cell. Whenever the network scores higher, what the
gap measures is that our fitting procedure did not find that map: a fixed ridge grid, a fixed number
of steps, and $8{,}192$ training states. It is not a statement that an affine readout cannot
represent what the network computes, because it demonstrably can.

This is why the $L^4$ differences of \PrimLfourHundredPostDiff{} to \ScPrimDiff{} are reported and
nothing is built on them. They are evidence about supervised recovery under a stated budget. A
reader who wants the claim ``the network's nonlinearity buys expressive power'' will not find it
here, and should not: the experiment as designed cannot supply it.

\paragraph{What the comparison does support.} The informative content is not any single cell but the
ordering across code families, because the protocol's asymmetries are held fixed while the code
varies. Both sides read the same state, both get one validation-selected global threshold, both are
scored on the same held-out supports. Under that fixed protocol the sign of the difference is set by
how the code was produced: positive and growing with width for $L^4$, exactly zero for $L^2$, and
negative by one to two and a half levels for an untrained code. A protocol that systematically
favoured the network could not produce a two-and-a-half-level loss on any arm.

So the comparison is usable as a property of the code and not as a claim about decoders. Read that
way it says something we did not expect: on an untrained code, a fitted affine probe recovers
supports the network's own thresholded readout misses by a wide margin, and on a code trained under
$L^4$ it does not. Training does not only improve the geometry the code presents
(Section~\ref{sec:res-primary}); it also brings the network's own readout close to the best one a
supervised fit of the same class will find, which is a form of the co-adaptation E8 isolates.

\paragraph{The check that would settle it, and why it is not here.} One experiment would decide
whether the residual $L^4$ gap is a fitting artefact: initialise the probe at the network's own
$(W_{\mathrm{out}},\theta)$ and refit. If the fitted probe then matches or exceeds the network, the
gap was optimisation; if the fit moves away and scores worse, the objective the probe minimises is
not aligned with exact support recovery, which would be a different and more interesting statement.
We did not run it, for the honest reason that we identified the confound after the campaigns were
complete and chose to report the ceiling rather than spend the remaining budget widening a result
the paper does not lean on. It is the first thing we would run next, and it costs a single pass over
saved weights, which are released.
